\documentclass[a4paper,12pt]{article}

\usepackage[utf8]{inputenc}
\usepackage[T5]{fontenc}
\usepackage[vietnamese]{babel}
\usepackage{graphicx}
\usepackage{titlesec}
\usepackage{setspace}
\usepackage{float}


\begin{document}

% ================== TRANG BÌA ==================
\thispagestyle{empty}

% ======= TRƯỜNG - KHOA =======
\begin{center}
{\large \textbf{TRƯỜNG ĐẠI HỌC SÀI GÒN}}\\[2pt]
{\large \textbf{KHOA CÔNG NGHỆ THÔNG TIN}}\\[40pt]

% ======= LOGO =======
\includegraphics[width=4cm]{logoITSGU.png}\\[40pt]

% ======= TÊN MÔN =======
{\Large \textbf{CÔNG NGHỆ PHẦN MỀM}}\\[14pt]

\rule{12cm}{0.4pt}\\[24pt]

% ======= TÊN BÀI TẬP LỚN =======
{\large \textbf{Bài Tập Lớn – Kiểm Thử Phần Mềm}}\\[8pt]

% ======= TÊN ĐỀ TÀI =======
{\normalsize Ứng dụng Đăng nhập \& Quản lý Sản phẩm}\\[4pt]
{\normalsize (Version 1.0)}\\[14pt]

\rule{12cm}{0.4pt}\\[24pt]

% ======= GVHD =======
\begin{tabular}{l l}
\textbf{GVHD:} & Từ Lăng Phiêu \\
\end{tabular}

% \textbf{GVHD:} \hspace{5pt} Từ Lăng Phiêu

\vfill

% ======= NGÀY THÁNG =======
{\normalsize TP. HỒ CHÍ MINH, THÁNG 11/2025}
\end{center}

\clearpage

% ================== MỤC LỤC ==================
\tableofcontents
\clearpage

% ================== NỘI DUNG CHÍNH ==================

\section{Giới thiệu về Dự án}

\subsection{Tổng quan}

Dự án \textbf{FLoginFE\_BE} là một ứng dụng web hoàn chỉnh bao gồm:
\begin{itemize}
    \item \textbf{Chức năng Login}: Hệ thống đăng nhập với validation đầy đủ
    \item \textbf{Chức năng Product}: Quản lý sản phẩm (CRUD operations)
    \item \textbf{Frontend}: React 18+
    \item \textbf{Backend}: Spring Boot 3.2+
    \item \textbf{Testing}: Phát triển theo phương pháp TDD
\end{itemize}

\subsection{Công nghệ sử dụng}

\subsubsection{Frontend}
\begin{itemize}
    \item React 18+ -- Framework JavaScript
    \item React Testing Library -- Testing cho React
    \item Jest -- Testing framework
    \item Axios -- HTTP client
    \item CSS3 -- Styling với animations
\end{itemize}

\subsubsection{Backend}
\begin{itemize}
    \item Spring Boot 3.2+ -- Framework Java
    \item Java 17+
    \item JUnit 5 -- Testing framework
    \item Mockito -- Mock framework
    \item Maven -- Build tool
    \item Spring Data JPA -- Database operations
\end{itemize}

\subsection{Cấu trúc dự án}

\textbf{FLoginFE\_BE/}
\begin{itemize}
    \item \textbf{frontend/} - React Application
    \begin{itemize}
        \item{src/}
        \begin{itemize}
            \item components/ -- Login, Product components
            \item services/ -- API services
            \item utils/ -- Validation utilities
            \item tests/ -- Test files
        \end{itemize}
        \item{package.json}
    \end{itemize}

    \item \textbf{backend/} - Spring Boot API
    \begin{itemize}
        \item{src/}
        \begin{itemize}
            \item  main/java/com/flogin/
            \begin{itemize}
                \item controller/ - AuthController, ProductController
                \item service/ - Business logic
                \item dto/ - Data Transfer Objects
                \item entity/ - Database entities
                \item repository/ - Data access
            \end{itemize}
            \item test/java/ - Test files
        \end{itemize}
        \item{pom.xl}
    \end{itemize}
\end{itemize}
\clearpage

\section{Câu 1: Phân tích và Thiết kế Test Cases (20 điểm)}

\subsection{Câu 1.1: Login - Phân tích và Test Scenarios (10 điểm)}

\subsubsection{Yêu cầu (5 điểm)}

Dựa trên chức năng đăng nhập (Login), hãy thực hiện:

a) Phân tích đầy đủ các yêu cầu chức năng của tính năng Login (2 điểm):

\begin{itemize}
    \item Validation rules cho username
    \item Validation rules cho password
    \item Authentication flow
    \item Error handling
\end{itemize}

b) Liệt kê và mô tả ít nhất 10 test scenarios cho Login bao gồm (2 điểm):

\begin{itemize}
    \item Happy path: Đăng nhập thành công
    \item Negative tests: Username/password rỗng, sai format
    \item Boundary tests: Độ dài min/max của username/password
    \item Edge cases: Ký tự đặc biệt, khoảng trắng
\end{itemize}

c) Phân loại test scenarios theo mức độ ưu tiên (Critical, High, Medium, Low) và giải thích (1 điểm)

\vspace{10pt}
\textbf{Validation Rules cho Login:}

\begin{itemize}
    \item Username: 3--50 ký tự, chỉ chứa a-z, A-Z, 0-9, ".", "-", "\_"
    \item Password: 6--100 ký tự, phải có cả chữ và số
\end{itemize}

\subsubsection{Thiết kế Test Cases chi tiết (5 điểm)}

Thiết kế 5 test cases quan trọng nhất cho Login theo template: 

\textit{Sinh viên cần thiết kế thêm 4 test cases tương tự}

\begin{table}[H]
\centering
\begin{tabular}{|l|l|}
\hline
\textbf{Test Case ID} & TC\_LOGIN\_001 \\ \hline
\textbf{Test Name} & Đăng nhập thành công với credentials hợp lệ \\ \hline
\textbf{Priority} & Critical \\ \hline
\textbf{Preconditions} & 
\begin{tabular}[c]{@{}l@{}}
- User account exists \\
- Application is running
\end{tabular} \\ \hline
\textbf{Test Steps} & 
\begin{tabular}[c]{@{}l@{}}
1. Navigate to login page \\
2. Enter valid username \\
3. Enter valid password \\
4. Click Login button
\end{tabular}
\\ \hline
\textbf{Test Data} &
\begin{tabular}[c]{@{}l@{}}
Username: testuser \\
Password: Test123
\end{tabular}
\\ \hline
\textbf{Expected Result} &
\begin{tabular}[c]{@{}l@{}}
- Success message displayed \\
- Token stored \\
- Redirect to dashboard
\end{tabular}
\\ \hline
\textbf{Actual Result} & (Để trống) \\ \hline
\textbf{Status} & Not Run \\ \hline
\end{tabular}
\caption{Template Test Case cho Login}
\end{table}


\subsection{Câu 1.2: Product – Phân tích và Test Scenarios (10 điểm)}

\subsubsection{Yêu cầu (5 điểm)}

Dựa trên chức năng quản lý sản phẩm (Product Management), hãy thực hiện:

a) Phân tích đầy đủ các yêu cầu chức năng của Product CRUD:

\begin{itemize}
    \item \textbf{Create}: Thêm sản phẩm mới
    \item \textbf{Read}: Xem danh sách/chi tiết sản phẩm
    \item \textbf{Update}: Cập nhật thông tin sản phẩm
    \item \textbf{Delete}: Xóa sản phẩm
\end{itemize}

b) Liệt kê và mô tả ít nhất 10 test scenarios cho Product, bao gồm:

\begin{itemize}
    \item \textbf{Happy path}: CRUD operations thành công
    \item \textbf{Negative tests}: Dữ liệu không hợp lệ
    \item \textbf{Boundary tests}: Giá trị min/max
    \item \textbf{Edge cases}: Sản phẩm trùng tên, xóa sản phẩm không tồn tại
\end{itemize}

c) Phân loại test scenarios theo mức độ ưu tiên và giải thích.

\textbf{Validation Rules cho Product}

\begin{itemize}
    \item Product Name: 3–100 ký tự, không được rỗng
    \item Price: $> 0$, $\le 999{,}999{,}999$
    \item Quantity: $\ge 0$, $\le 99{,}999$
    \item Description: $\le 500$ ký tự
    \item Category: Phải thuộc danh sách categories có sẵn
\end{itemize}


\subsubsection{Thiết kế Test Cases chi tiết (5 điểm)}
Thiết kế 5 test cases quan trọng nhất cho Product Management:

\begin{table}[H]
\centering
\begin{tabular}{|l|l|}
\hline
\textbf{Test Case ID} & TC\_PRODUCT\_001 \\ \hline
\textbf{Test Name} & Tạo sản phẩm mới thành công \\ \hline
\textbf{Priority} & Critical \\ \hline
\textbf{Preconditions} & 
\begin{tabular}[c]{@{}l@{}}
- User đã đăng nhập \\
- User có quyền tạo sản phẩm
\end{tabular} \\ \hline
\textbf{Test Steps} & 
\begin{tabular}[c]{@{}l@{}}
1. Navigate to Product page \\
2. Click "Add New Product" \\
3. Enter product information \\
4. Click Save button
\end{tabular}
\\ \hline
\textbf{Test Data} &
\begin{tabular}[c]{@{}l@{}}
Name: Laptop Dell \\
Price: 15000000 \\
Quantity: 10 \\
Category: Electronics
\end{tabular}
\\ \hline
\textbf{Expected Result} &
\begin{tabular}[c]{@{}l@{}}
- Product created successfully \\
- Success message displayed \\
- Product appears in list
\end{tabular}
\\ \hline
\textbf{Actual Result} & (Để trống) \\ \hline
\textbf{Status} & Not Run \\ \hline
\end{tabular}
\caption{Template Test Case cho Login}
\end{table}

\textbf{Sinh viên cần thiết kế thêm 4 test cases tương tự cho Update, Delete, Read operations}

\section{Câu 2: Unit Testing và Test-Driven Development (20 điểm)}

\subsection{Câu 2.1: Login - Unit Tests Frontend và Backend (10 điểm)}

\subsubsection{3.1.1 Frontend Unit Tests - Validation Login (5 điểm)}

Áp dụng TDD để phát triển unit tests cho validation module của Login.

\textbf{Yêu cầu:}

a) Viết unit tests cho \texttt{validateUsername()} (2 điểm):

\begin{itemize}
    \item Test username rỗng
    \item Test username quá ngắn/dài
    \item Test ký tự đặc biệt không hợp lệ
    \item Test username hợp lệ
\end{itemize}

b) Viết unit tests cho \texttt{validatePassword()} (2 điểm):

\begin{itemize}
    \item Test password rỗng
    \item Test password quá ngắn/dài
    \item Test password không có chữ hoặc số
    \item Test password hợp lệ
\end{itemize}

c) Coverage $\geq 90\%$ cho validation module (1 điểm)

\textbf{Ví dụ Frontend Unit Test:}

\begin{figure}[H]
    \centering
    \includegraphics[width=0.8\textwidth]{download.jpeg}
    \caption{Template Test Case cho Login}
    \label{fig:login_case}
\end{figure}

\subsubsection{Backend Unit Tests - Login Service (5 điểm)}

Viết unit tests cho \texttt{AuthService} của backend:

\textbf{Yêu cầu:}

a) Test method \texttt{authenticate()} với các scenarios (3 điểm):

\begin{itemize}
    \item Login thành công
    \item Login với username không tồn tại
    \item Login với password sai
    \item Validation errors
\end{itemize}

b) Test validation methods riêng lẻ (1 điểm)

c) Coverage $\geq 85\%$ cho AuthService (1 điểm)

\textbf{Ví dụ Backend Unit Test:}

\begin{figure}[H]
    \centering
    \includegraphics[width=0.8\textwidth]{download.jpeg}
    \caption{AuthServiceTest.java – ví dụ minh họa Backend Unit Test}
    \label{fig:auth_service_test}
\end{figure}


\subsection{Câu 2.2: Product - Unit Tests Frontend và Backend (10 điểm)}

\subsubsection{Frontend Unit Tests - Product Validation (5 điểm)}

Áp dụng TDD cho validation của Product:

\textbf{Yêu cầu:}

a) Viết unit tests cho \texttt{validateProduct()} (3 điểm):

\begin{itemize}
    \item Test product name validation
    \item Test price validation (boundary tests)
    \item Test quantity validation
    \item Test description length
    \item Test category validation
\end{itemize}

b) Viết tests cho Product form component (1 điểm)

c) Coverage $\geq 90\%$ (1 điểm)

\textbf{Ví dụ Product Validation Test:}

\begin{figure}[H]
    \centering
    \includegraphics[width=0.8\textwidth]{download.jpeg}
    \caption{AuthServiceTest.java – ví dụ minh họa Backend Unit Test}
    \label{fig:auth_service_test}
\end{figure}

\subsubsection{ Backend Unit Tests - Product Service (5 điểm)}

Viết unit tests cho \texttt{ProductService}:

\textbf{Yêu cầu:}

a) Test CRUD operations (4 điểm):

\begin{itemize}
    \item Test createProduct()
    \item Test getProduct()
    \item Test updateProduct()
    \item Test deleteProduct()
    \item Test getAll() với pagination
\end{itemize}

b) Coverage $\geq 85\%$ cho ProductService (1 điểm)

\textbf{Ví dụ Product Service Test:}

\begin{figure}[H]
    \centering
    \includegraphics[width=0.8\textwidth]{download.jpeg}
    \caption{AuthServiceTest.java – ví dụ minh họa Backend Unit Test}
    \label{fig:auth_service_test}
\end{figure}

\section{Câu 3: Integration Testing (20 điểm)}

\subsection{Câu 3.1: Login - Integration Testing (10 điểm)}

\subsubsection{4.1.1 Frontend Component Integration (5 điểm)}

Test tích hợp Login component với API service:

\textbf{Yêu cầu:}

a) Test rendering và user interactions (2 điểm)

b) Test form submission và API calls (2 điểm)

c) Test error handling và success messages (1 điểm)

\textbf{Ví dụ Component Integration Test:}

\begin{figure}[H]
    \centering
    \includegraphics[width=0.8\textwidth]{download.jpeg}
    \caption{AuthServiceTest.java – ví dụ minh họa Backend Unit Test}
    \label{fig:auth_service_test}
\end{figure}

\subsubsection{4.1.2 Backend API Integration (5 điểm)}

Test API endpoints của Login với MockMvc:

\textbf{Yêu cầu:}

a) Test POST /api/auth/login endpoint (3 điểm)

b) Test response structure và status codes (1 điểm)

c) Test CORS và headers (1 điểm)

\textbf{Ví dụ API Integration Test:}

\begin{figure}[H]
    \centering
    \includegraphics[width=0.8\textwidth]{download.jpeg}
    \caption{AuthServiceTest.java – ví dụ minh họa Backend Unit Test}
    \label{fig:auth_service_test}
\end{figure}

\subsection{Câu 3.2: Product - Integration Testing (10 điểm)}

\subsubsection{4.2.1 Frontend Component Integration (5 điểm)}

Test tích hợp Product components:

\textbf{Yêu cầu:}

a) Test ProductList component với API (2 điểm)

b) Test ProductForm component (create/edit) (2 điểm)

c) Test ProductDetail component (1 điểm)

\textbf{Ví dụ Product Component Test:}

\begin{figure}[H]
    \centering
    \includegraphics[width=0.8\textwidth]{download.jpeg}
    \caption{AuthServiceTest.java – ví dụ minh họa Backend Unit Test}
    \label{fig:auth_service_test}
\end{figure}

\subsubsection{4.2.2 Backend API Integration (5 điểm)}

Test các API endpoints của Product:

\textbf{Yêu cầu:}

a) Test POST /api/products (Create) (1 điểm)

b) Test GET /api/products (Read all) (1 điểm)

c) Test GET /api/products/\{id\} (Read one) (1 điểm)

d) Test PUT /api/products/\{id\} (Update) (1 điểm)

e) Test DELETE /api/products/\{id\} (Delete) (1 điểm)

\textbf{Ví dụ Product API Test:}

\begin{figure}[H]
    \centering
    \includegraphics[width=0.8\textwidth]{download.jpeg}
    \caption{AuthServiceTest.java – ví dụ minh họa Backend Unit Test}
    \label{fig:auth_service_test}
\end{figure}

\section{Câu 4: Mock Testing (10 điểm)}

\subsection{Câu 4.1: Login - Mock Testing (5 điểm)}

\subsubsection{5.1.1 Frontend Mocking (2.5 điểm)}

Mock external dependencies cho Login component:

\textbf{Yêu cầu:}

a) Mock \texttt{authService.loginUser()} (1 điểm)

b) Test với mocked successful/failed responses (1 điểm)

c) Verify mock calls (0.5 điểm)

\textbf{Ví dụ Frontend Mock Test:}

\begin{figure}[H]
    \centering
    \includegraphics[width=0.8\textwidth]{download.jpeg}
    \caption{AuthServiceTest.java – ví dụ minh họa Backend Unit Test}
    \label{fig:auth_service_test}
\end{figure}

\subsubsection{5.1.2 Backend Mocking (2.5 điểm)}

Mock dependencies trong Backend tests:

\textbf{Yêu cầu:}

a) Mock \texttt{AuthService} với \texttt{@MockBean} (1 điểm)

b) Test controller với mocked service (1 điểm)

c) Verify mock interactions (0.5 điểm)

\textbf{Ví dụ Backend Mock Test:}

\begin{figure}[H]
    \centering
    \includegraphics[width=0.8\textwidth]{download.jpeg}
    \caption{AuthServiceTest.java – ví dụ minh họa Backend Unit Test}
    \label{fig:auth_service_test}
\end{figure}

\subsection{Câu 4.2: Product - Mock Testing (5 điểm)}

\subsubsection{5.2.1 Frontend Mocking (2.5 điểm)}

Mock \texttt{ProductService} trong component tests:

\textbf{Yêu cầu:}

a) Mock CRUD operations (1.5 điểm)

b) Test success và failure scenarios (0.5 điểm)

c) Verify all mock calls (0.5 điểm)

\textbf{Ví dụ Product Mock Test:}

\begin{figure}[H]
    \centering
    \includegraphics[width=0.8\textwidth]{download.jpeg}
    \caption{AuthServiceTest.java – ví dụ minh họa Backend Unit Test}
    \label{fig:auth_service_test}
\end{figure}

\subsubsection{5.2.2 Backend Mocking (2.5 điểm)}

Mock \texttt{ProductRepository} trong service tests:

\textbf{Yêu cầu:}

a) Mock \texttt{ProductRepository} (1 điểm)

b) Test service layer với mocked repository (1 điểm)

c) Verify repository interactions (0.5 điểm)

\textbf{Ví dụ Repository Mock Test:}
\begin{figure}[H]
    \centering
    \includegraphics[width=0.8\textwidth]{download.jpeg}
    \caption{AuthServiceTest.java – ví dụ minh họa Backend Unit Test}
    \label{fig:auth_service_test}
\end{figure}

\section{Câu 5: Automation Testing và CI/CD (10 điểm)}

\subsection{Câu 5.1: Login - E2E Automation Testing (5 điểm)}

\subsubsection{6.1.1 Setup và Configuration (1 điểm)}

\textbf{Yêu cầu:}

\begin{itemize}
    \item Cài đặt Cypress hoặc Selenium
    \item Cấu hình test environment
    \item Setup Page Object Model
\end{itemize}

\subsubsection{6.1.2 E2E Test Scenarios cho Login (2.5 điểm)}

Viết automated tests cho toàn bộ login flow:

\textbf{Yêu cầu:}

a) Test complete login flow (1 điểm)

b) Test validation messages (0.5 điểm)

c) Test success/error flows (0.5 điểm)

d) Test UI elements interactions (0.5 điểm)

\textbf{Ví dụ Cypress E2E Test:}

\begin{figure}[H]
    \centering
    \includegraphics[width=0.8\textwidth]{download.jpeg}
    \caption{AuthServiceTest.java – ví dụ minh họa Backend Unit Test}
    \label{fig:auth_service_test}
\end{figure}

\subsubsection{6.1.3 CI/CD Integration cho Login Tests (1.5 điểm)}

\textbf{Yêu cầu:}

\begin{itemize}
    \item Tạo GitHub Actions workflow
    \item Run login tests automatically
    \item Generate test reports
\end{itemize}

\textbf{Ví dụ CI/CD Config:}

\begin{figure}[H]
    \centering
    \includegraphics[width=0.8\textwidth]{download.jpeg}
    \caption{AuthServiceTest.java – ví dụ minh họa Backend Unit Test}
    \label{fig:auth_service_test}
\end{figure}

\subsection{Câu 5.2: Product - E2E Automation Testing (5 điểm)}

\subsubsection{6.2.1 Setup Page Object Model (1 điểm)}

Implement POM cho Product pages:

\textbf{Ví dụ Page Object:}

\begin{figure}[H]
    \centering
    \includegraphics[width=0.8\textwidth]{download.jpeg}
    \caption{AuthServiceTest.java – ví dụ minh họa Backend Unit Test}
    \label{fig:auth_service_test}
\end{figure}

\subsubsection{6.2.2 E2E Test Scenarios cho Product (2.5 điểm)}

Viết automated tests cho CRUD operations:

\textbf{Yêu cầu:}

a) Test Create product flow (0.5 điểm)

b) Test Read/List products (0.5 điểm)

c) Test Update product (0.5 điểm)

d) Test Delete product (0.5 điểm)

e) Test Search/Filter functionality (0.5 điểm)

\textbf{Ví dụ Product E2E Test:}

\begin{figure}[H]
    \centering
    \includegraphics[width=0.8\textwidth]{download.jpeg}
    \caption{AuthServiceTest.java – ví dụ minh họa Backend Unit Test}
    \label{fig:auth_service_test}
\end{figure}

\textbf{Ví dụ Product E2E Test (tiếp):}

\begin{figure}[H]
    \centering
    \includegraphics[width=0.8\textwidth]{download.jpeg}
    \caption{AuthServiceTest.java – ví dụ minh họa Backend Unit Test}
    \label{fig:auth_service_test}
\end{figure}

\subsubsection{6.2.3 CI/CD Integration (1.5 điểm)}

Setup complete CI/CD pipeline:

\textbf{Ví dụ Complete CI/CD:}

\begin{figure}[H]
    \centering
    \includegraphics[width=0.8\textwidth]{download.jpeg}
    \caption{AuthServiceTest.java – ví dụ minh họa Backend Unit Test}
    \label{fig:auth_service_test}
\end{figure}

\section{Phần Mở Rộng (Bonus 20 điểm)}

\subsection{7.1 Performance Testing (10 điểm)}

\subsubsection{7.1.1 Yêu cầu}

a) Setup JMeter hoặc k6 cho performance testing (2 điểm)

b) Viết performance tests cho Login API (3 điểm):

\begin{itemize}
    \item Load test: 100, 500, 1000 concurrent users
    \item Stress test: Tìm breaking point
    \item Response time analysis
\end{itemize}

c) Viết performance tests cho Product API (3 điểm)

d) Phân tích kết quả và đưa ra recommendations (2 điểm)

\subsection{7.2 Security Testing (10 điểm)}

\subsubsection{7.2.1 Yêu cầu}

a) Test common vulnerabilities (5 điểm):

\begin{itemize}
    \item SQL Injection
    \item XSS (Cross-Site Scripting)
    \item CSRF (Cross-Site Request Forgery)
    \item Authentication bypass attempts
\end{itemize}

b) Test input validation và sanitization (3 điểm)

c) Security best practices implementation (2 điểm):

\begin{itemize}
    \item Password hashing
    \item HTTPS enforcement
    \item CORS configuration
    \item Security headers
\end{itemize}

\section{Tiêu chí Đánh giá}

\subsection{8.1 Bảng Phân bổ Điểm}

\begin{table}[H]
\centering
\begin{tabular}{|l|c|c|c|}
\hline
\textbf{Nội dung} & \textbf{Login} & \textbf{Product} & \textbf{Tổng} \\ \hline
Câu 1: Phân tích và Test Cases & 10 & 10 & 20 \\ \hline
Câu 2: Unit Testing và TDD & 10 & 10 & 20 \\ \hline
Câu 3: Integration Testing & 10 & 10 & 20 \\ \hline
Câu 4: Mock Testing & 5 & 5 & 10 \\ \hline
Câu 5: Automation và CI/CD & 5 & 5 & 10 \\ \hline
\textbf{Tổng điểm bắt buộc} & 40 & 40 & 80 \\ \hline
Phần Mở Rộng (Bonus) & \multicolumn{2}{c|}{} & +20 \\ \hline
\textbf{Điểm tối đa} & \multicolumn{2}{c|}{} & 100 \\ \hline
\end{tabular}
\caption{Bảng phân bổ điểm chi tiết}
\end{table}

\subsection{8.2 Tiêu chí Chất lượng}

\subsubsection{8.2.1 Code Quality (30\%)}

\begin{itemize}
    \item Clean code principles
    \item Proper test structure (AAA pattern)
    \item Meaningful test names
    \item Test coverage $\geq 80\%$
    \item All tests pass
\end{itemize}

\subsubsection{8.2.2 Documentation (20\%)}

\begin{itemize}
    \item Test cases đầy đủ và rõ ràng
    \item Screenshots và evidences
    \item Test reports
    \item README với hướng dẫn chạy tests
\end{itemize}

\subsubsection{8.2.3 Completeness (30\%)}

\begin{itemize}
    \item Hoàn thành tất cả yêu cầu bắt buộc
    \item Đầy đủ test cases cho cả Login và Product
    \item CI/CD pipeline hoạt động
\end{itemize}

\subsubsection{8.2.4Best Practices (20\%)}

\begin{itemize}
    \item Áp dụng TDD đúng cách
    \item Proper mocking strategy
    \item Good test data management
    \item Good test data management
    \item Automation best practices
\end{itemize}

\section{Hướng dẫn Nộp bài}

\subsection{9.1 Format Nộp bài}

\subsubsection{9.1.1 Source Code}

\begin{itemize}
    \item Git repository (GitHub/GitLab) – Public hoặc add giảng viên
    \item Commit history rõ ràng
    \item README.md đầy đủ
    \item Có .gitignore phù hợp
\end{itemize}

\subsubsection{9.1.2 Báo cáo}

\begin{itemize}
    \item File PDF (tối đa 50 trang)
    \item Bao gồm:
    \begin{itemize}
        \item Giới thiệu project
        \item Test cases chi tiết (Login + Product)
        \item Test execution results
        \item Screenshots và evidences
        \item Coverage reports
        \item CI/CD pipeline documentation
        \item Kết luận
    \end{itemize}
\end{itemize}

\subsection{9.2 Thời hạn}

\begin{itemize}
    \item \textbf{Deadline: [11/11/2025]}
    \item Không chấp nhận nộp trễ
\end{itemize}

\subsection{9.3 Demo (Optional)}

\begin{itemize}
    \item Video demo 10–15 phút
    \item Demo chạy tests
    \item Giải thích CI/CD pipeline
    \item Q\&A với giảng viên
\end{itemize}

\section*{Tài liệu}

\begin{enumerate}
    \item React Documentation, \textbf{https://react.dev}, Testing Library Documentation
    \item Spring Boot Documentation, \textbf{https://spring.io/projects/spring-boot}, Spring Testing Guide
    \item Jest Documentation, \textbf{https://jestjs.io/}, JavaScript Testing Framework
    \item JUnit 5 User Guide, \textbf{https://junit.org/junit5/}, Testing Framework for Java
    \item Mockito Framework, \textbf{https://site.mockito.org/}, Mocking Framework for Java
    \item Cypress Documentation, \textbf{https://www.cypress.io/}, End-to-End Testing Framework
    \item Test-Driven Development: By Example, \textit{Kent Beck}, Addison-Wesley Professional
\end{enumerate}


\end{document}
